\section{Spanning-tree normal form}

Choose a root $r\in V(X)$ and a spanning tree $T\subseteq X$. Orient each
tree edge away from $r$.

\begin{proposition}[Tree normalization]
Every receipt assignment is gauge equivalent to one whose receipt is the
identity on every dart of $T$.
\end{proposition}

\begin{proof}
Set $g(r)=e$. If a tree edge $a:u\to v$ points from a parent to a child,
define recursively
\[
g(v)=\alpha(a)^{-1}g(u).
\]
Then
\[
\alpha^g(a)=g(u)^{-1}\alpha(a)g(v)=e.
\]
The reverse tree dart is also trivial. The tree contains a unique root path
to every vertex, so the recursion is well defined.
\end{proof}

Let
\[
\beta=|E|-|V|+1
\]
be the cycle rank of $X$. There are exactly $\beta$ undirected cotree edges.
After tree normalization, choosing one orientation of each cotree edge gives
a tuple
\[
(c_1,\ldots,c_{\beta})\in R^{\beta}.
\]
Each $c_j$ is the receipt of the corresponding fundamental circuit, because
the tree portions of that circuit have trivial receipt.

The fundamental group of a connected graph is free of rank $\beta$
\cite{serre-trees}. Therefore the tuple freely determines a homomorphism
\[
\rho_{\alpha}:\pi_1(X,r)\longrightarrow R.
\]
The spanning tree chooses free generators; it does not create the represented
path-receipt law.
